% META: {"id": "1NSI-TAB-M4", "chapitre": "1NSI-TABLES", "type_objet": "methode", "capacites": ["P-TABLE-04"], "capacites_codes": ["C4"], "methodes": ["M4"], "mode_creation": "ex_nihilo", "sources_inspiration": [], "status": "generated"}
\begin{fichemethode}{M4}{Fusionner deux tables par une clé commune}
\textbf{Quand l'utiliser ?} Quand l'information est répartie sur deux fichiers et qu'une
question porte sur les deux à la fois : les notes d'un côté, les coordonnées de l'autre,
reliées par un identifiant.

\textbf{Pas à pas :}
\begin{enumerate}
  \item \textbf{Identifie la clé commune} : la colonne présente dans les deux tables et qui
    désigne la même chose. Vérifie qu'elle est bien \textbf{unique} dans au moins l'une des
    deux — sinon la fusion multipliera les lignes.
  \item \textbf{Indexe la seconde table} par cette clé, dans un dictionnaire
    \lstinline!{valeur_de_cle: ligne}!. C'est l'étape qui fait tout : sans elle, on
    reparcourt la seconde table pour chaque ligne de la première.
  \item \textbf{Parcours la première table} et, pour chaque ligne, va chercher la ligne
    correspondante dans l'index.
  \item \textbf{Décide du sort des lignes sans correspondance} : les écarter, ou les garder
    avec des champs vides. Ce choix change le résultat — écris-le noir sur blanc.
  \item \textbf{Projette} : construis la ligne fusionnée en ne gardant que les colonnes
    utiles, et sans écraser deux colonnes de même nom.
\end{enumerate}

\textbf{Exemple :}
\begin{python}
def fusionner(gauche, droite, cle):
    """Joint deux tables sur `cle`.

    Les lignes sans correspondance sont ecartees.
    """
    index = {ligne[cle]: ligne for ligne in droite}
    resultat = []
    for ligne in gauche:
        appariee = index.get(ligne[cle])
        if appariee is not None:
            resultat.append({**ligne, **appariee})
    return resultat
\end{python}

% BEGIN-VERIFY
% notes = [
%     {"id": "e1", "points": 17},
%     {"id": "e2", "points": 9},
%     {"id": "e9", "points": 12},
% ]
% eleves = [
%     {"id": "e1", "nom": "Ines"},
%     {"id": "e2", "nom": "Karim"},
%     {"id": "e3", "nom": "Lou"},
% ]
% def fusionner(gauche, droite, cle):
%     index = {ligne[cle]: ligne for ligne in droite}
%     resultat = []
%     for ligne in gauche:
%         appariee = index.get(ligne[cle])
%         if appariee is not None:
%             resultat.append({**ligne, **appariee})
%     return resultat
% jointe = fusionner(notes, eleves, "id")
% assert [l["nom"] for l in jointe] == ["Ines", "Karim"]
% assert jointe[0] == {"id": "e1", "points": 17, "nom": "Ines"}
% # e9 n'a pas d'eleve, e3 n'a pas de note : les deux disparaissent
% assert len(jointe) == 2 < min(len(notes), len(eleves)) + 1
% assert "e9" not in [l["id"] for l in jointe]
% assert "Lou" not in [l["nom"] for l in jointe]
% # une cle dupliquee a droite : l'index n'en garde qu'une, silencieusement
% doublons = [{"id": "e1", "nom": "Ines"}, {"id": "e1", "nom": "Autre"}]
% index = {l["id"]: l for l in doublons}
% assert len(index) == 1 and index["e1"]["nom"] == "Autre"
% # colonnes homonymes : la droite ecrase la gauche
% assert {**{"x": 1}, **{"x": 2}} == {"x": 2}
% END-VERIFY

\textbf{Pièges :}
\begin{itemize}
  \item \textbf{Une clé dupliquée dans la table indexée.} Le dictionnaire n'en garde que la
    dernière, sans le moindre message. Contrôle d'abord :
    \lstinline{len(index) == len(droite)}.
  \item \textbf{Ne pas décider du sort des lignes orphelines.} Ici elles sont écartées, donc
    la table fusionnée est plus courte que les deux. Si on attendait toutes les lignes de
    gauche, le résultat est silencieusement amputé.
  \item \textbf{Deux colonnes de même nom dans les deux tables.} Avec
    \lstinline!{**gauche, **droite}!, celle de droite écrase celle de gauche. Renomme avant
    de fusionner si les deux comptent.
  \item Comparer des clés de types différents : \lstinline{"12"} et \lstinline{12} ne
    s'apparient jamais. C'est encore la conversion oubliée à l'import.
  \item \textbf{Choisir comme clé une colonne qui n'identifie rien.} \lstinline{genre} ou
    \lstinline{classe} se répètent : la fusion apparie alors des lignes qui n'ont rien à voir
    et le résultat paraît plausible. Une clé doit désigner \textbf{une} ligne et une seule.
\end{itemize}

\textbf{Vérifier :} compare \lstinline{len(jointe)} au nombre de clés réellement communes.
Si la fusion est plus \textbf{longue} que les deux tables, une clé est dupliquée ; si elle
est plus courte que prévu, des lignes n'ont pas trouvé de correspondance.

\textbf{S'entraîner :} exercices C4 du chapitre.
\end{fichemethode}
